\begin{theorem}[Exterior replacement with prescribed winding]%
    \label{thm:peps_torus_complement_path_replacement}
    \lean{TNLean.PEPS.exists_complementWalk_of_integerExteriorCollarGraph_connected,
        TNLean.PEPS.exists_complementBoundaryLoops_of_integerExteriorCollarGraph_connected}
    \leanok
    \uses{thm:peps_integer_exterior_collar_paths,
        thm:peps_torus_region_boundary_lift,
        thm:peps_torus_exterior_path_approximation}
    Let both torus periods be at least three. Suppose a region has an
    integer unit-step lift, its safe exterior graph is connected, and
    its exterior collar projects outside the actual torus region.
    Every route which enters through one crossing bond, follows an
    interior path, and leaves through another crossing bond has a
    complementary replacement with the same two seam crossing numbers.

    In particular, fix rooted spanning trees in the region and its
    complement and a reference crossing bond. For each crossing bond
    there is a loop based at the complementary root, wholly within
    the induced complement, with the winding of the corresponding
    two-tree boundary route. The loops are chosen independently of
    all group-valued closure labels. This is an auxiliary geometric
    construction for \cite[proof of Theorem~6.9]{Schuch2010PEPS}.
\end{theorem}
\begin{proof}\leanok
    Lift the two exterior crossing endpoints. They belong to the
    integer band, and the deck-coordinate difference equals the
    crossing numbers of the interior route. A safe graph walk gives
    a continuous collar path between them. Projection exclusion and
    lattice-path approximation give a native complementary walk with
    that same deck-coordinate difference. Prepend the complementary
    tree path from the root to the first crossing and append the
    tree path from the second crossing back to the root. Crossing
    numbers add under concatenation, proving the loop assertion.
\end{proof}

\begin{theorem}[Complementary boundary loops of a simply connected region]%
    \label{thm:peps_torus_simply_connected_boundary_loops}
    \lean{TNLean.PEPS.exists_complementBoundaryLoops_of_isSimplyConnected}\leanok
    \uses{thm:peps_torus_complement_path_replacement,
        thm:peps_integer_boundary_contour_connected,
        thm:peps_torus_integer_cell_lift}
    Let both torus periods be at least three, and suppose the actual
    closed-cell realization of $R$ is simply connected. Choose rooted
    spanning trees in $R$ and its induced complement and fix a
    reference crossing bond. There is a complementary root loop for
    every crossing bond with the same winding as its two-tree boundary
    route. These loops are chosen before all group-valued closure labels.
    No exterior path, collar connectivity, or winding identity is an
    additional premise. This supplies the exterior geometric construction
    in \cite[proof of Theorem~6.9]{Schuch2010PEPS}.
\end{theorem}
\begin{proof}\leanok
    Choose one covering lift whose planar integer-cell union is simply
    connected and whose exterior collar projects outside the torus
    region. Its safe exterior graph is connected by the boundary contour
    theorem. Apply the preceding replacement construction to the
    chosen two-tree routes.
\end{proof}
